\section{仿真实验}
    本小节将详细介绍异构卫星网络路由算法的仿真实验环境、参数设置及对比分析结果。实验旨在验证本文所提多域协同的分层路由机制（记为H-IGA），在多层异构、高动态及非均匀业务负载场景下的性能优势，特别是针对链路拥塞的缓解能力。
    \subsection{仿真环境与参数设置}
    1、仿真场景
    
    仿真实验基于 Python 语言开发，采用 STK (Systems Tool Kit) 软件生成高精度的卫星轨道星历与链路可见性数据。仿真场景构建了与\ref{sec: STK参数配置}一致的异构天地一体化网络场景，包含100 颗低轨探测卫星、225 颗低轨通信卫星、8 颗中轨卫星、3 颗地球同步轨道卫星、48 个全球分布的地面站，全网共计 384 个节点。仿真利用 STK 导出的真实星历数据，以 10秒 为时间步长，动态计算节点位置与链路可见性。为了模拟真实的异构网络特性，不同层级间的链路带宽与通信距离限制采用了非均匀配置，如表\ref{tab: 异构网络链路参数配置}所示。
    
    \begin{table}[htbp]
        \caption{异构网络链路参数配置}
        \centering
        \begin{tabular}{lccc}
        \toprule
        \textbf{链路类型} & \textbf{带宽 (Mbps)} & \textbf{传播时延 (ms)} & \textbf{基础丢包率 (\%)} \\
        \midrule
        Ground $\leftrightarrow$ LEO/Detect & 100 & 2.4 $\sim$ 13.2 & $0.1 \sim 0.5$ \\
        Ground $\leftrightarrow$ MEO/GEO & 100 & 50.0 $\sim$ 130.0 & $0.1 \sim 0.5$ \\
        MEO $\leftrightarrow$ MEO/GEO & 100 & 60.0 $\sim$ 150.0 & $0.01 \sim 0.1$ \\
        LEO/Detect $\leftrightarrow$ MEO & 60 & 20.0 $\sim$ 45.5 & $0.01 \sim 0.1$ \\
        LEO/Detect $\leftrightarrow$ GEO & 60 & 117.0 $\sim$ 135.0 & $0.01 \sim 0.1$ \\
        Detect $\leftrightarrow$ LEO & 30 & 1.0 $\sim$ 24.0 & $0.01 \sim 0.1$ \\
        LEO/Detect $\leftrightarrow$ LEO/Detect & 30 & 1.4 $\sim$ 15.0 & $0.01 \sim 0.1$ \\
        \bottomrule
        \end{tabular}
        \label{tab: 异构网络链路参数配置}
    \end{table}


    
    2、链路物理模型设置
    
    为了真实反映网络负载对传输质量的影响，仿真采用了基于 BPR (Bureau of Public Roads) 模型计算动态时延，并结合过载丢包惩罚机制。链路的总时延 $D_{total}$ 由传播时延 $D_{prop}$ 和动态排队时延 $D_{dyn}$ 组成：
    $$ D_{total} = D_{prop} \times (1 + \alpha \cdot u^\beta) $$
    其中 $u$ 为链路利用率 ($u = \text{UsedBW} / \text{Capacity}$)，$\alpha$ 和 $\beta$ 为拥塞敏感系数，仿真中分别设置为 1.0 和 6.0，以模拟高负载下时延的剧烈增长。动态时延最大惩罚倍数限制为 500。
    
    丢包率 $L$ 同样基于负载率动态计算，模拟拥塞导致的队列溢出：
    $$
    L(u) =
    \begin{cases}
      0.1\% + u \times 0.1\% & u \le 1.0 \\
      0.2\% + (u-1.0) \times 80\% & u > 1.0 \text{ (过载丢包)}
    \end{cases}
    $$
    当链路过载时，丢包率随负载率迅速上升，最高限制为 99\%。
    
    3、业务流量模型
    
    实验构建热点区域与背景流量混合的非均匀业务场景，以测试算法在局部高负载下的性能。对于业务分布，80\%的业务源节点集中在北美西海岸 (Lat: $30^{\circ}\sim40^{\circ}$, Lon: $-125^{\circ}\sim-115^{\circ}$)，宿节点集中在东亚区域 (Lat: $30^{\circ}\sim40^{\circ}$, Lon: $130^{\circ}\sim145^{\circ}$)，这种跨洋长距离传输强迫流量汇聚于有限的 MEO 节点，制造激烈的跨层资源竞争；20\%的业务在全球范围内随机选取源宿节点对，用于模拟背景通信负载。实验定义了四种不同 QoS 需求的业务流，如表\ref{tab: 多级 QoS 业务类型定义}所示。
    \begin{table}[htbp]
      \caption{多级 QoS 业务类型定义}
      \centering
      \begin{tabular}{lcccc}
        \toprule
        \textbf{业务类型} & \textbf{带宽 (Mbps)} & \textbf{时延容限 (ms)} & \textbf{丢包率容限 (\%)} & \textbf{优先级}  \\
        \midrule
        遥感业务 (Remote Sensing) & 20  & 500 & 5.0\% & 吞吐量优先 \\
        视频业务 (Video Live) & 10  & 150 &  2.0\% & 时延优先 \\
        话音业务 (Voice VoIP) & 1   & 100 & 2.0\% & 时延优先 \\
        测控业务 (Telemetry Control) & 0.1 & 80 & 1.0\% & 高可靠/低时延 \\
        \bottomrule
      \end{tabular}
      \label{tab: 多级 QoS 业务类型定义}
    \end{table}
    
    \subsubsection{算法参数设置}
    为了验证本文提出的分层拥塞感知算法 (H-IGA) 的性能，实验将其与以下两种典型算法进行对比：
    \begin{enumerate}
        \item \textbf{Dijkstra 算法}：经典的单源最短路径算法。在仿真中，该算法直接在扁平化的全网拓扑上运行，以链路的静态传播时延作为权重，不具备分层架构与负载感知能力，作为性能底线（Baseline）。
        \item \textbf{分层标准遗传算法 (SGA)}：同样采用多域分层路由架构，但在域内路由阶段使用标准遗传算法。与 H-IGA 不同，SGA 不包含方向引导初始化机制，且其适应度函数仅考虑路径时延与丢包的简单倒数，不具备非线性的拥塞惩罚项。其种群规模设置为 30，最大迭代次数为 30。
    \end{enumerate}
    
    H-IGA 的核心进化参数设置如表 4-3 所示。
    
    \begin{table}[htbp]
      \caption{H-IGA 算法参数设置}
      \centering
      \begin{tabular}{lc}
        \toprule
        \textbf{参数名称} & \textbf{数值} \\
        \midrule
        种群规模 (Population Size) & 25 \\
        最大迭代次数 (Max Iterations) & 25 \\
        交叉概率 ($P_c$) & 0.8 \\
        变异概率 ($P_m$) & 0.2 \\
        方向引导概率 ($P_{guide}$) & 0.7 \\
        拥塞敏感度 ($\gamma$) & 2.0 \\
        拥塞感知权重 ($\lambda$) & 1.0 \\
        \bottomrule
      \end{tabular}
      \label{tab:iga_params}
    \end{table}
    
    针对不同业务的 QoS 需求，H-IGA 在计算适应度时采用动态权重系数。对于时延优先型业务（测控、视频、话音），时延权重 $\alpha=0.8$，丢包权重 $\beta=0.2$；对于吞吐量优先型业务（遥感），时延权重 $\alpha=0.2$，丢包权重 $\beta=0.8$。

    

    